\section{The marked decision spectrum}\label{sec:v8-spectrum}
We first fix the object whose finite presentations are to be compared.
This requires specifying both the decision tasks and the information
available to the decoder. These specifications cannot be recovered
from a dimension or a cardinality alone.

\subsection{Experiments, tasks, and signatures}
An experiment is a standard Borel causal instrument, with an initial law
and transition--observation kernels indexed by a parameter $\theta$.
Actions, when present, are chosen before the corresponding transition.
The same observation-based policy is used at every value of $\theta$.
A \emph{mark} $d$ specifies a horizon, decision times, action spaces, a
loss in $[0,1]$, and a risk functional. The latter is either expectation
under a declared prior or the supremum over a declared parameter set,
taken after selecting one decoder. A minimax decoder does not know the
unknown parameter. A mark may instead specify a fixed Bayesian
prediction problem with an integrable squared loss; assertions involving
total variation below apply only after a uniform loss bound is supplied.
When a loss concerns a latent target, the marked experiment includes that
target, and a simulation must preserve its joint law with the observations.
Equality of report laws alone does not suffice for such a mark.

An encoder $Q$ is an initialized causal transducer
\begin{equation}
 I_t\sim q_t(\,\cdot\mid I_{t-1},Y_t,A_{t-1}),
 \qquad I_t\in\{1,\ldots,S\}.                         \label{eq:v8-register}
\end{equation}
After an update, a decoder sees only $(t,I_t)$ and any explicitly
declared current side input. It does not see the discarded observation,
an output transcript, old actions, or a data-dependent random seed.
All persistent decoder state belongs to $I_t$. Fresh randomization is
allowed. A public seed, when allowed, is independent of the experiment;
a data-dependent position or modification of that seed is charged.
A terminal reconstruction kernel may have a continuous output even when
the physical statistic has a finite alphabet.

We use the resource signature
\begin{equation}
 b=(S,A,L,P,p,W,\tau;\mathsf R).                       \label{eq:v8-signature}
\end{equation}
Here $S$ counts persistent states; $A$ counts physical terminal labels;
$L$ bounds raw symbols read per acquisition; $P$ bounds fixed description
bits; $p$ is numerical precision; $W$ is transient work space; $\tau$ is
update time; and $\mathsf R$ records the permitted independent
randomization. A coordinate $\infty$ means that no bound is asserted.
For a precision constraint, the order is the order on permitted
implementations, not the numerical order on an accuracy tolerance.
The \emph{cardinality signature} constrains only $S$ (and $A$ when
explicitly stated). In particular, exact real constants in a fixed
program are free only under that signature. All theorems below state
which signature they use.

For an encoder admissible at $b$, let $r_E(Q,d)$ be the infimal value of
the declared risk functional over its allowed decoders. Let $\beta_E(d)$ be the infimal risk using the full
observation history under the same decision timing. The encoder is
selected before the mark; its decoder may depend on the mark. For a set
$\mathcal D$ of marks define
\begin{align}
 \mathscr A_E(b)&=\{u\in[0,1]^{\mathcal D}:\text{some admissible }Q
                  \text{ has }r_E(Q,d)\le u_d\ (d\in\mathcal D)\},
                                                               \label{eq:v8-riskset}\\
 \Theta_E(b;\mathcal D)&=\inf_Q\sup_{d\in\mathcal D}
                         \{r_E(Q,d)-\beta_E(d)\}.       \label{eq:v8-spectrum}
\end{align}
Infima need not be attained. We use the displayed attainable sets, not
their product-topology closure: for an infinite task family, taking that
closure before a uniform task supremum can change the value. The common
encoder and the order of quantifiers in \eqref{eq:v8-spectrum} are part
of the invariant.
The collection of attainable sets, baselines, and resource signatures is
the \emph{marked decision spectrum}, denoted $\boldsymbol\Theta(E)$.
This is the meaning of the symbol $\Theta$ in the present paper. It is a
profile, not an additional physical parameter. Unbounded squared marks
are treated separately without restricting the vectors to $[0,1]$.

\begin{example}[The encoder cannot follow the task]\label{ex:v8-quantifiers}
Let $X=(X_1,X_2)$ consist of two independent fair bits. A terminal
statistic has two values. The two marks ask, respectively, for $X_1$ and
$X_2$, with zero--one loss. Each mark separately admits risk zero, but
one common encoder has worst risk at least $1/4$. Indeed, for each fixed
independent coding seed and each encoder label, the two optimal answers
form a word in $\{0,1\}^2$. Two reconstruction words have average Hamming
distance at least $1/2$ from the four equally likely input words: two
words may be exact, and each remaining word is at distance at least one.
Randomization does not lower the average of the two risks. Hence their
maximum is at least $1/4$. Conversely, with a public fair coin, reveal
one of the two coordinates and guess the other. Thus the value is $1/4$
when that coin is permitted, whereas interchanging $\inf_Q$ and the task
supremum gives zero. No such interchange is made in this paper.
\end{example}

\subsection{Typed simulations and invariance}
A certified simulation $\mathsf S:E\Rightarrow F$ consists of initialized
causal kernels, a monotone resource transformer $\rho_{\mathsf S}$, and
an error bound $\epsilon_{\mathsf S}(T)$. The kernels do not know $\theta$
or the latent variables. For every target-side admissible policy and
every parameter, the simulated marked transcript through horizon $T$
has total variation distance at most $\epsilon_{\mathsf S}(T)$ from its
law in $F$. Moreover, composing the simulation with any target encoder
of signature $b$ produces a source encoder of signature at most
$\rho_{\mathsf S}(b)$. These are uniform requirements on one simulator,
not permission to choose a different simulator for each task or policy.
When marks have different horizons, use the corresponding error for each
coordinate; a common $\epsilon$ below denotes their supremum.

Two certified simulators represent the same morphism when they have the
same transformer and error certificate and induce the same marked
execution laws under every admissible continuation. Internal register
names are immaterial. We use standard Borel kernels modulo this
execution equivalence. The exact morphisms have error zero. A
$K$-state cardinality simulator has the transformer $S\mapsto KS$;
its own output is not a second free memory register.

\begin{theorem}[Spectrum transport and resource equivalence]
\label{thm:v8-spectrum}
For bounded marks and signatures as in \eqref{eq:v8-signature}, certified
simulations have identities and associative composition up to execution
equivalence. If $\mathsf S:E\Rightarrow F$ and
$\mathsf T:F\Rightarrow G$, the composite has transformer
$\rho_{\mathsf S}\circ\rho_{\mathsf T}$ and error at most
$\min\{1,\epsilon_{\mathsf S}+\epsilon_{\mathsf T}\}$.
For every target-admissible $Q$ there is a source-admissible $Q'$ with
\begin{equation}
 r_E(Q',d)\le r_F(Q,d)+\epsilon_{\mathsf S},\qquad
 \beta_E(d)\le\beta_F(d)+\epsilon_{\mathsf S}.          \label{eq:v8-rawtransport}
\end{equation}
If a reverse simulation $\mathsf R:F\Rightarrow E$ is also supplied,
then
\begin{equation}
 \Theta_E(\rho_{\mathsf S}(b);\mathcal D)
 \le\Theta_F(b;\mathcal D)+\epsilon_{\mathsf S}+\epsilon_{\mathsf R}.
                                                               \label{eq:v8-excesstransport}
\end{equation}
In particular, mutually exact finite-state simulations give equal ideal
baselines and comparable spectra under the two budget dilations.
Whenever a positive power exponent of either cardinality profile exists,
the other has the same exponent. These assertions also hold coordinatewise
for the attainable risk sets, before taking the task supremum.
\end{theorem}
\begin{proof}
Compose the kernels on the product of the two internal states. Fix one
policy for $G$, regard its composition with $\mathsf T$ as a policy and
continuation for $F$, and apply the certificate for $\mathsf S$.
Data processing followed by the certificate for $\mathsf T$ gives the
sum of the two errors. The composition of transformers is forced by the
same execution. Associativity follows from iterated kernel integration
and the canonical reassociation of the triple product of registers;
execution equivalence is preserved under composition because its
quantifier includes all continuations. The identity kernel, identity
transformer, and zero error give identities. This establishes a category
for exact simulations and the corresponding error calculus for approximate
ones; it is not an assertion of a new general randomization theorem.

Run $Q$ on the simulated reports and use its decoder. Every loss in
$[0,1]$ changes in expectation by at most the variation error. Taking
infima over decoders proves the first inequality in
\eqref{eq:v8-rawtransport}; an arbitrarily good full-history policy proves
the second. The reverse simulator implies
$\beta_F(d)-\beta_E(d)\le\epsilon_{\mathsf R}$. Subtract the baselines,
take the common task supremum, and only then take the infimum over $Q$.
This proves \eqref{eq:v8-excesstransport}. No compactness or attainment of
an optimal encoder is used.

For exact cardinality equivalence with factors $K_1,K_2$, the inequalities
are $\Theta_E(K_1M)\le\Theta_F(M)$ and
$\Theta_F(K_2M)\le\Theta_E(M)$. Monotonicity in $M$ and
$\log(K_iM)/\log M\to1$ sandwich any existing limit of
$-\log\Theta(M)/\log M$. The same reasoning applies to a single bounded
squared prediction mark; multiplication by its loss bound changes no
positive power exponent.
\end{proof}

The theorem is an invariance statement for a \emph{typed} equivalence.
Ordinary statistical equivalence without a resource certificate need not
preserve a causal implementation budget. Conversely, retaining every
emitted label in an external transcript changes the signature and thus
changes the spectrum. Products of independent observations and products
of memories are different operations: the former concern the experiment,
the latter the implementation.

\begin{proposition}[Forgetting causal and resource structure]
\label{prop:v8-forgetful}
Fix a finite horizon and restrict to uncontrolled instruments and
parameter-decision marks. Sending an instrument to its full transcript
experiment, and an exact causal simulator to its transcript kernel,
defines a functor to parameterized statistical experiments with
parameter-independent kernels modulo equality on every parameter law.
It preserves identities and composition. It forgets the resource
certificate; the resulting statistical equivalence alone does not
supply an equality of causal-register spectra.
\end{proposition}
\begin{proof}
Iterating a simulator's initialized kernels at a fixed input transcript
gives a parameter-independent Markov kernel from source transcripts to
target transcripts. Iterated integration identifies the transcript
kernel of a composite with the composite of these kernels. Execution
equivalence implies equality on every source parameter law, so this
assignment is well defined on the stated classes. The identity gives
the identity transcript kernel. A general inverse transcript kernel
need not be causal, and a causal inverse may store a nonfinite state;
neither datum is reconstructed by the forgetful operation.
\end{proof}

\subsection{The terminal decision-complete slice}
In this subsection $E=\{P_\theta:\theta\in K\}$ has standard Borel
observation space $X$, $K$ is compact metric, and $\theta\mapsto P_\theta$
is continuous in total variation. A terminal encoder is any Markov kernel
$Q:X\to[S]$. Only its physical alphabet $A\le S$ is constrained; transient
computation and the reconstruction kernel are unrestricted. Let
$\mathcal D_{\rm all}$ consist of all finitely supported priors on $K$
and all finite-action losses in $[0,1]$. For these families we use
$\delta(F,E)=\inf_R\sup_{\theta\in K}\|RF_\theta-P_\theta\|_{\rm TV}$,
where $R$ is a parameter-independent reconstruction kernel and
$\|P-Q\|_{\rm TV}=\sup_B|P(B)-Q(B)|$.

\begin{theorem}[Decision completeness at a fixed common encoder]
\label{thm:v8-duality}
For every such $Q$,
\begin{equation}
 \sup_{d\in\mathcal D_{\rm all}}\{r_E(Q,d)-\beta_E(d)\}
       =\delta(QE,E).
                                                               \label{eq:v8-duality} 
\end{equation}
Consequently
\begin{equation}
 \Theta_E(S;\mathcal D_{\rm all})
       =\inf_{Q:X\to[S]}\delta(QE,E).                 \label{eq:v8-terminalprofile}
\end{equation}
This is the Blackwell--Le Cam randomization criterion in a compact
setting, with the common finite encoder kept outside the decision
supremum. In particular, no new priority is claimed for the
unconstrained criterion \cite{Blackwell1953,LeCam1964,Torgersen1991}.
\end{theorem}
\begin{proof}
A reconstruction $R$ with variation error $\epsilon$ transfers each
bounded decision rule, proving that the left side is at most the right
side. First suppose that both $K$ and $X$ are finite. Put $F=QE$.
Finite-dimensional convex minimax applied to the row-stochastic matrices
$R$ and the total-variation dual gives
\[
 \delta(F,E)=\sup_{\pi,h}\left[
 \sum_{\theta,x}\pi_\theta h_{\theta x}P_\theta(x)
 -\sup_R\sum_{\theta,i,x}\pi_\theta F_\theta(i)
                         R(i,x)h_{\theta x}\right],
\]
where $\pi$ is a prior and $0\le h_{\theta x}\le1$. To justify convexity
of the dual domain, write its variables as
$u_{\theta x}=\pi_\theta h_{\theta x}$ with
$0\le u_{\theta x}\le\pi_\theta$ and $\sum\pi_\theta=1$.
Regard $h$ as a reward with action set $X$. The first term is the reward
of the identity decision in $E$, bounded above by its optimal reward;
the second is the optimal reward in $F$. Thus every dual expression is
at most one of the optimal reward differences. Changing reward to loss
proves equality.

Here is the approximation needed for general $X,K$. A dense sequence
$(\theta_j)$ and positive weights with sum one give
$m=\sum_jw_jP_{\theta_j}$ dominating every $P_\theta$; domination follows
by total-variation continuity. The set of densities
$\{dP_\theta/dm:\theta\in K\}$ is compact in $L^1(m)$. Given $\eta>0$,
choose a finite $L^1$ net and approximate its elements by simple
functions on a common finite partition $\mathcal C$ of $X$. Conditional
expectation onto $\mathcal C$ is an $L^1$ contraction, so the cell
reconstruction $R_{\mathcal C}(\cdot\mid C)=m(\cdot\mid C)$ satisfies
\[
 \sup_\theta\|R_{\mathcal C}\mathcal C P_\theta-P_\theta\|_{\rm TV}
 \le\eta
\]
after choosing the net and approximations sufficiently accurately.
Zero-mass cells are irrelevant. Hence
$\delta(F,E)\le\delta(F,\mathcal CE)+\eta$.
Now choose a finite parameter net on which both $F$ and $\mathcal CE$
vary by at most $\eta$ in total variation. The supremum reconstruction
error on $K$ differs from that on the net by at most $2\eta$, uniformly
in the reconstruction. Apply the finite result to this net. Its priors
are allowed priors on $K$, and observing $X$ is at least as informative
as observing $\mathcal C(X)$, so the resulting decision differences for
$E$ dominate those for $\mathcal CE$. It follows that the left side of
\eqref{eq:v8-duality} is at least $\delta(F,E)-3\eta$. Let $\eta\downarrow0$.
Equation \eqref{eq:v8-terminalprofile} now follows by taking the infimum
over $Q$, with no exchange of quantifiers.
\end{proof}
